\subsection{无穷小变换}

“无限小”变换的概念可以追溯到无穷小演算的开端；众所周知，笛卡儿通过假设“在无限小之中”每个平面运动都可以等同于一个旋转而发现了瞬时旋转中心；十八世纪解析力学的精心构建完全建立在类似的想法之上。1851 年，西尔维斯特在设法构造线性群

GL(3, C) 或其某些子群的不变量时，给这些矩阵中出现的参数 \(z_{j}\) 以形如 \(\alpha_{j}dt\) 的“无限小”增量，并表达函数 \(f((z_{j}))\) 的不变性，即写下方程 \(f((z_{j} + \alpha_{j}dt)) = f((z_{j}))\)；这给了他一个关于 f 的、带偏导数的线性方程 Xf = 0，其中

\[
X f = \sum_ {j} \alpha_ {j} \frac {\partial f}{\partial z _ {j}},\tag{25.1}
\]

X 因而是一个微分算子，即“一个沿有向参数 \(\alpha_{j}\) 的方向的导数”（{[}304{]}，vol.~3，pp.~326 与 327）；西尔维斯特似乎感觉到其中有一个适用范围相当广的一般原理，但似乎没有再回到这个问题上来。稍后，凯莱（{[}58{]}，v. II, pp.~164-178）对 SL(2, C) 在该群的某些表示中的不变量以同样的方式处理，并证明它满足两个一阶偏微分方程 Xf = 0, Yf = 0，其中 X 与 Y 如上那样从“无限小”变换出发而得到

\[
\left( \begin{array}{c c} 0 & 0 \\ d t & 0 \end{array} \right) \text {   and   } \left( \begin{array}{c c} 0 & d t \\ 0 & 0 \end{array} \right).
\]

用现代的语言，这可以由 X 与 Y 生成李代数 \(\mathfrak{sl}(2,\mathbf{C})\) 这一事实来解释；此外，凯莱还明确地算出了换位积 XY-YX，并证明它也来自一个“无限小”变换。

在若尔当 1868 年关于运动群的论文{[}174 b{]}中，“无限小变换”的概念自始至终被使用，但纯粹是从几何的观点出发。毫无疑问，正是从他那里出现了由一个无限小变换“生成”的单参数群的思想：对若尔当而言，它是通过“适当地重复”该无限小变换而得到的变换全体（loc. cit.，p.~243）。克莱因与李在他们 1871 年的论文中使用了同样的说法“重复的无限小变换”（{[}182{]}，v. I, pp.~424-459），但上下文表明，他们用这个词指的是一个微分方程组的积分。如果他们所考虑的单参数群由变换 \(x' = f(x, y, t)\)，\(y' = g(x, y, t)\) 组成，则相应的“无限小变换”由

\[
d x = p (x, y) d t, d y = q (x, y) d t
\]

其中 \(p(x,y)=\frac{\partial f}{\partial t}(x,y,t_{0}),q(x,y)=\frac{\partial g}{\partial t}(x,y,t_{0})\)，这里 \(t_{0}\) 对应于群的恒等变换。由于克莱因与李显式地知道函数 f 与 g，他们毫不困难地验证了函数

\[
t \mapsto f (x, y, t) \text {   and   } t \mapsto g (x, y, t)
\]

以参数形式给出了微分方程

\[
q (\xi , \eta) d \xi = p (\xi , \eta) d \eta
\]

的经过点 \((x,y)\) 的积分曲线，但对这一点他们没有给出任何一般的理由；此外，在他们的论文的其余部分，他们也没有再用到这一事实。

